\chapter{Удиртгал}
\label{ch:introduction}

\section{Үндэслэл, ач холбогдол}

Өнөө үед гар утас нь хүмүүсийн өдөр тутмын амьдралын салшгүй хэсэг болж, хүн бүр түүнийгээ байнга биедээ авч явдаг болсон. Үүнтэй зэрэгцэн тэмдэглэл хөтлөх, төлөвлөгөө боловсруулах үйл ажиллагааг дижитал хэлбэрт шилжүүлэх хэрэгцээ нэмэгдэж байгаа ч уламжлалт цаасан тэмдэглэл нь мартагдах, гээгдэх зэрэг сул талтай.

Иймд гар утасны аппликейшн хэлбэрт шилжүүлснээр тэмдэглэлийг ямар ч үед хөтлөх, мэдээллийг аюулгүй хадгалах, сануулга ашиглан ажлаа цаг тухайд нь биелүүлэх боломж бүрдэнэ.

Судалгааны ажлын зорилго нь хиймэл оюуныг ашиглан хэрэглэгчийн тэмдэглэл, төлөвлөгөөний агуулгад тохирсон стикерийг автоматаар үүсгэж, түүнийг календар хэсэгт харуулах боломжтой, хялбар ойлгомжтой интерфэйстэй дижитал тэмдэглэлийн аппликэйшн хөгжүүлэхэд оршино.

Энэхүү зорилгын хүрээнд системийн шаардлага, бүтэц, үйл ажиллагааг тодорхойлж, Flutter ашиглан хэрэглэгчийн харагдах хэсгийг хөгжүүлэн, Node.js/Express болон MongoDB ашиглан сервер талын хэсэг, өгөгдлийн санг хэрэгжүүлэв. Мөн OpenRouter API ашиглан AI sticker generation болон notification функцийг нэгтгэн, ``AI Notebook'' системийг боловсруулж хэрэгжүүлэв.

Мөн судалгааны ажлын ач холбогдол нь өдөр тутмын тэмдэглэл хөтлөх энгийн хэрэглээг хиймэл оюуны тусламжтай контентын баяжуулалттай хослуулж, хэрэглэгчийн оролцоо, тогтвортой хэрэглээ, мэдээллийн хүртээмжийг сайжруулсан систем санал болгосонд оршино. Ийм хандлага нь цаашид боловсрол, хувийн бүтээмж, сэтгэлзүйн өөрийгөө ажиглах зэрэг салбарт өргөжих боломжтой.

\section{Зорилго, зорилт}

\subsection{Судалгааны зорилго}

Энэхүү судалгааны ажлын гол зорилго нь харагдах хэсэг (frontend) болон Node.js/Express дээрх сервер талын хэсэг (backend)-д тулгуурласан, хэрэглэгчийн баталгаажуулалт, тэмдэглэлийн менежмент, хиймэл оюун ухаанаар стикер үүсгэх (AI sticker generation) боломж бүхий ухаалаг тэмдэглэлийн гар утасны апп зохион бүтээж, хөгжүүлэн, түүний ажиллагаа, үр нөлөөг туршилтаар үнэлэхэд оршино.

\subsection{Судалгааны зорилтууд}

Дээрх зорилгод хүрэхийн тулд дараах зорилтуудыг дэвшүүлсэн:

\begin{enumerate}
    \item Системийн шаардлага, бүтэц, үйл ажиллагааг тодорхойлох.
    \item Flutter ашиглан хэрэглэгчдэд зориулсан хялбар, ойлгомжтой апп хөгжүүлэх.
    \item Node.js/Express болон MongoDB ашиглан backend server, өгөгдлийн сантай холбох.
    \item OpenRouter API ашиглан sticker generation болон notification функцийг хэрэгжүүлэх.
    \item Системийн бүрэн үйл ажиллагаанд туршилт хийх.
\end{enumerate}

\section{Судалгааны ажлын бүтэц}

Энэхүү төгсөлтийн судалгааны ажил дараах бүтэцтэйгээр зохион байгуулагдсан.

\begin{enumerate}
    \item Нэгдүгээр бүлэгт судалгааны үндэслэл, зорилго, зорилт тодорхойлсон.
    \item Хоёрдугаар бүлэгт сэдэвтэй холбоотой онолын суурь, хэрэглэгдэж буй технологийн өнөөгийн түвшинг тоймлосон.
    \item Гуравдугаар бүлэгт системийн арга зүй, өгөгдлийн урсгалын түвшинчилсэн задлан шинжилгээ, дарааллын диаграммуудыг (sequence diagrams) тайлбарласан.
    \item Дөрөвдүгээр бүлэгт хэрэгжилтийн үр дүн, интеграцийн туршилт, баталгаажуулалтын дүнг нэгтгэн харуулсан.
    \item Тавдугаар бүлэгт судалгааны ерөнхий дүгнэлт, хязгаарлалт, цаашдын хөгжүүлэлтийн чиглэлийг тодорхойлсон.
\end{enumerate}
